% CHECKPOINT36: actual pure abelianization and the derived-transporter strengthening.
% Independent geometric argument; no checkpoint35 source is changed.
\statement{A.125}{Proved}{The elliptic transporter lies in the pure derived subgroup}
Retain the actual groups \(\mathcal G,\mathcal P\), element
\(\Pi_0\), and unique pure transporters \(p_w,h_w\) of A.124.
Write \(P=\widehat{\mathcal P}\), \(P'=\overline{[P,P]}\), and
\(\rho_0=\operatorname{Ad}(\Pi_0)|_P\). Then
\[
 \mathcal P_{\mathrm{ab}}\simeq\mathbb Z^4,\qquad
 P_{\mathrm{ab}}\simeq\widehat{\mathbb Z}^{\,4},\qquad
 (P_{\mathrm{ab}})^{\rho_0}=0.                               \tag{A.125.1}
\]
For every profinite commutator word \(w\) satisfying the original
pentagon,
\[
 p_w\in P',\qquad
 (P')^{\rho_0}=1,\qquad H^1(C_3,P';\rho_0)=\{*\}.              \tag{A.125.2}
\]
The map \(w\mapsto p_w\) is continuous with values in the actual
closed subgroup \(P'\).
For the actual elliptic involution \(\Theta\) of A.122,
\[
 h_w\in P',\qquad z_w=h_w^{-1}\in P'.                         \tag{A.125.3}
\]
Thus A.124's constant section obstruction has a specified refinement
in \(H^1(C_2,P';\Theta)\), constant for all these \(w\).
Its image is the class in \(H^1(C_2,P;\Theta)\) of A.124.
No injectivity of this last map, and no triviality of either
specified class, is asserted.

\begin{proof}
\emph{1. Integral homology of the actual configuration space.}
By the marked geometric identification in A.124,
\(\mathcal P=\pi_1(\mathcal X)\), where
\[
 \begin{gathered}
 M=T^2\times T^2,\qquad
 \mathcal X=M\setminus(D_1\cup D_2\cup D_3),\\
 D_1=\{a=0\},\qquad D_2=\{b=0\},\qquad D_3=\{a=b\}.
 \end{gathered}
\]
Here \(T^2=\mathbb R^2/\mathbb Z^2\); thus \(M\) has real
dimension four. The three oriented codimension-two subtori
meet only in \(o=(0,0)\).

Remove \(o\) from \(M\). The relative homology of
\((M,M\setminus\{o\})\) vanishes in degrees below four,
as follows by excision to a four-ball and its punctured ball.
The long exact sequence of the triple therefore identifies
\[
 H_i(M,\mathcal X;\mathbb Z)
 =H_i(M\setminus\{o\},\mathcal X;\mathbb Z)
 \quad (i\leq2).
\]
In \(M\setminus\{o\}\), the three punctured tori
\(D_j\setminus\{o\}\) are disjoint closed connected submanifolds
with oriented real normal two-plane bundles. Choose disjoint
tubular neighborhoods, with radius allowed to decrease near
the omitted point. Excision reduces the relative calculation
to their disk bundles modulo their punctured disk bundles.
In degrees one and two, the relative cellular calculation has
one oriented normal disk in degree two at each base point;
paths in a connected base identify these generators with the
same sign. There is no relative cell below degree two.
Equivalently, this is the low-degree oriented Thom calculation.
It gives
\[
 H_2(M,\mathcal X;\mathbb Z)=\mathbb Z^3,\qquad
 H_1(M,\mathcal X;\mathbb Z)=0.                               \tag{A.125.4}
\]
The map \(H_2(M;\mathbb Z)\to\mathbb Z^3\) is the vector
of oriented intersection numbers with the three \(D_j\).

This map is onto, as can be checked by three actual tori.
For a fixed nonzero \(b_0\), the torus \(a\mapsto(a,b_0)\)
has intersection vector \((1,0,1)\).
For a fixed nonzero \(a_0\), the torus \(b\mapsto(a_0,b)\)
has vector \((0,1,1)\).
Finally choose nonzero \(\alpha,\beta\in\mathbb R/\mathbb Z\)
and use
\[
 (t,u)\longmapsto
       \bigl(a=(t,\alpha),\ b=(\beta,u)\bigr).
\]
It avoids \(D_1,D_2\) and meets \(D_3\) once, transversely,
at \(t=\beta,\ u=\alpha\). Its vector is \((0,0,\pm1)\).
These three vectors form an integral basis. The relative
homology sequence and (A.125.4) consequently prove that
inclusion induces the actual isomorphism
\[
 H_1(\mathcal X;\mathbb Z)\xrightarrow{\ \sim\ }
 H_1(M;\mathbb Z)=\mathbb Z^4.                               \tag{A.125.5}
\]
No unspecified presentation relation has been discarded.

\emph{2. The cyclic action and full profinite abelianization.}
A three-cycle on normalized configurations is
\((a,b)\mapsto(b-a,-a)\), or its inverse, according to the
choice of cyclic direction. On each of the two torus-coordinate
blocks its matrix is therefore
\[
 C=\begin{pmatrix}-1&1\\-1&0\end{pmatrix}
 \quad\text{or}\quad C^{-1}.
\]
Both have characteristic polynomial \(t^2+t+1\), and both
\(C-I\) and \(C^{-1}-I\) have determinant three.
The action of the actual lift \(\Pi_0\) on pure abelianization
is this permutation action: any change of a lift is an inner
pure action and disappears on abelianization.
The whole rank-four determinant is nine.
Equation (A.125.5) is equivariant, so these are the actual
integral blocks for \(\rho_0\).

The universal finite-abelian quotient property gives
\[
 P_{\mathrm{ab}}
   =\widehat{\mathcal P_{\mathrm{ab}}}
   =\widehat{\mathbb Z}^{\,4}.
\]
Multiplication by three, and hence by nine, is injective
on \(\widehat{\mathbb Z}\). The adjugate matrix identity
therefore shows that \(\rho_0-I\) has zero kernel on
\(P_{\mathrm{ab}}\), despite not being onto at three.
This proves (A.125.1); invertibility at every prime is
neither needed nor claimed.

\emph{3. The exact derived transporter.}
All arguments in the three evaluations defining \(\Pi_w\)
are pure:
\[
 (YX^{-1}Y^{-1},K),\qquad (X,K),\qquad
 (\sigma^2,\tau^2).
\]
Since \(w\) is a profinite commutator word, each evaluation
lies in \(P'\). This subgroup is normal in
\(\widehat{\mathcal G}\), because \(P\) is normal and its
closed derived subgroup is characteristic. Thus the literal
word \(\Pi_w=Lw_1Yw_2^{-1}w_3\) gives
\[
 \Pi_w\Pi_0^{-1}\in P'.
\]
Using the unique transporter of A.124, this difference is
\[
 p_w\Pi_0p_w^{-1}\Pi_0^{-1}
          =p_w\rho_0(p_w)^{-1}.
\]
Its image in \(P_{\mathrm{ab}}\) is zero, so
\((1-\rho_0)([p_w])=0\). By (A.125.1), \([p_w]=0\).
Therefore \(p_w\in P'\). The continuous selection of A.124
already takes its values in this closed subgroup, and is
continuous for its actual subspace topology.

The fixed subgroup \((P')^{\rho_0}\) is contained in
\(C_{\widehat{\mathcal G}}(\Pi_0)\cap P=1\).
For a continuous \(C_3\)-cocycle \(a\in P'\), its equation is
\[
 a\,\rho_0(a)\,\rho_0^2(a)=1.
\]
Hence \(a\Pi_0\) has cube one. It has the same nonzero
abelian label and the same permutation as \(\Pi_0\), so
A.124 gives a unique pure \(p\) with
\(a\Pi_0=p\Pi_0p^{-1}\).
The preceding abelian argument applies to this equality
as well, because \(a\in P'\), and gives \(p\in P'\).
Thus \(a=p\rho_0(p)^{-1}\), an actual coboundary
(put \(g=p^{-1}\) in \(g^{-1}\rho_0(g)\)).
This proves the entire \(H^1\) assertion in (A.125.2).

\emph{4. The refined involution obstruction.}
The actual section preserves \(P\), hence also \(P'\).
Put
\[
 F_b=\phi(\sigma^2,\tau^2),\quad
 \Phi=\phi(X,K),\quad Q=\phi(YX^{-1}Y^{-1},K).
\]
All three words lie in \(P'\), since \(\phi\) is a
commutator word and all their arguments are pure.
The original section formulas are
\[
 \Theta(\sigma)=F_b\sigma^{-1}F_b^{-1},\qquad
 \Theta(\tau)=\tau^{-1},\qquad
 \Theta(Y)=K^{-1}QY\Phi^{-1}.
\]
Consequently, modulo \(P'\),
\[
 \Theta(\Pi_0)=\Theta(LY)
       \equiv\tau\sigma K^{-1}Y
       =\tau^{-1}\sigma^{-1}Y=\Pi_0,                          \tag{A.125.6}
\]
where the exact equality uses \(K=\sigma\tau^2\sigma\).
The unique pure \(h_1\) with
\(\Theta(\Pi_0)=h_1\Pi_0h_1^{-1}\) therefore has
\((1-\rho_0)([h_1])=0\). Thus \(h_1\in P'\).

For general \(w\), the uniqueness formula of A.124 is
\[
 h_w=\Theta(p_w)h_1p_w^{-1}.
\]
All its factors lie in \(P'\). Thus \(h_w,z_w\in P'\),
and \(z_w\Theta(z_w)=1\) is a cocycle in that coefficient
group itself. The formula
\(z_w=p_wz_1\Theta(p_w)^{-1}\) proves constancy of the
refined class using correcting elements in \(P'\).
Functoriality gives its stated image in \(H^1(C_2,P;\Theta)\).
Triviality of the refined class is exactly existence of a
\(P'\)-conjugate of \(\Pi_w\) fixed by \(\Theta\).
The proof does not identify this with triviality after
enlarging the coefficient group, or evaluate the actual class.
\end{proof}
